\section*{第 \ref{ch:b2:precambrian-earth} 章——前寒武纪地球}

\begin{solution}{exo:b2:precambrian-earth:1}
(a) \omterm{def:b2:precambrian-earth:eons}{冥古宙} $4567 - 4031 = \qty{536}{Ma}$；\omterm{def:b2:precambrian-earth:eons}{太古宙} $4031 - 2500 = \qty{1531}{Ma}$；
\omterm{def:b2:precambrian-earth:eons}{元古宙} $2500 - 538.8 = \qty{1961}{Ma}$。(b) $(4567 - 538.8)/4567 = 0.882$：约
八分之七。(c) 锆石为 $4567 - 4404 = \qty{163}{Ma}$；最古老的岩石为
$4567 - 4030 = \qty{537}{Ma}$。
\end{solution}

\begin{solution}{exo:b2:precambrian-earth:2}
(a) 因子 $e^{\lambda t}$（其中 $\lambda = \ln2/T$，$t = \qty{3}{Ga}$）分别为 1.593
（\ce{^{238}U}）、19.18（\ce{^{235}U}）、1.160（\ce{^{232}Th}）和 5.292（\ce{^{40}K}）。各项
贡献分别为 3.053、1.573、2.424 和 \qty{4.329}{pW}，合计 \qty{11.38}{pW}，
而今日为 \qty{4.906}{pW}。(b) $11.38/4.906 = 2.32$。(c) 今日以钍-232 居首；当时
则以钾-40 居首。
\end{solution}

\begin{solution}{exo:b2:precambrian-earth:3}
(a) $1616 + 18.0\times(28 - 32.2) = \qty{1540}{\celsius}$。(b) $32.2 + (1500 - 1616)/18.0 =
\qty{25.8}{\percent}$。(c) 全岩分析并不等同于液体成分：
\omterm{def:b2:precambrian-earth:greenstone}{科马提岩}熔流中可能含有从液体中堆积起来的橄榄石晶体，它们会增加 MgO。其
\qty{38}{\percent} 超出了最富镁的\omterm{def:b2:precambrian-earth:greenstone}{科马提岩}液体的 \qty{32}{\percent}，并且
也超出了该直线所依据的实验。
\end{solution}

\begin{solution}{exo:b2:precambrian-earth:4}
最古老的岩石，\qty{4.03}{Ga}；最早的叠层石，\qty{3.43}{Ga}；皮尔巴拉穹窿结束，
\qty{3.1}{Ga}；\omterm{def:b2:precambrian-earth:goe}{大氧化事件}，约 2.4 至 \qty{2.1}{Ga}；最后一批常见\omterm{def:b2:precambrian-earth:greenstone}{科马提岩}，
\qty{1.9}{Ga}；氧气第二次上升，约 \qty{0.8}{Ga}；马里诺冰期结束，
\qty{0.635}{Ga}；最早的埃迪卡拉纪化石，约 \qty{0.57}{Ga}。
\end{solution}

\begin{solution}{exo:b2:precambrian-earth:5}
(a) 吸收的能量为 $\tfrac14\times1361\times0.706 = \qty{240.2}{W.m^{-2}}$，而 $\sigma\times288^4 =
\qty{390.1}{W.m^{-2}}$，故 $\varepsilon = 240.2/390.1 = 0.616$。(b) $\tfrac14\times1361\times0.4
= \qty{136.1}{W.m^{-2}}$, $T^4 = 136.1/(0.616\times\num{5.670e-8})$, $T = \qty{249.8}{K}$.
(c) 在 \qty{275.5}{K} 时，反照率为 $0.6 - 0.306\times12.5/20 = 0.4088$：$R = 340.25\times0.5912 =
\qty{201.2}{W.m^{-2}}$，且 $E = 0.616\,\sigma\times275.5^4 = \qty{201.2}{W.m^{-2}}$。(d) $R' =
340.25\times0.306/20 = \qty{5.21}{W.m^{-2}.K^{-1}}$, $E' = 4\varepsilon\sigma T^3 =
\qty{2.92}{W.m^{-2}.K^{-1}}$：$R' > E'$，故平衡态不稳定
（见\cref{prop:b2:precambrian-earth:ebm}）。
\end{solution}

\begin{solution}{exo:b2:precambrian-earth:6}
(a) $1/(1 + 0.4\times2.4/4.567) = 0.826$；$1/(1 + 0.4\times0.65/4.567) = 0.946$。(b) 暖态
需要至少今日太阳光的 0.932：\qty{2.4}{Ga} 前不可能，
\qty{0.65}{Ga} 前则可能。(c) $0.826/0.946 = 0.873$：暗了 \qty{13}{\percent}，接近
所引的 \qty{10}{\percent}，该数值来自另一种太阳模型。
\end{solution}

\begin{solution}{exo:b2:precambrian-earth:7}
(a) 斯图特冰期 $716.9 - 658.9 = \qty{58.0}{Ma}$；马里诺冰期 $644.5 - 635.4 = \qty{9.1}{Ma}$；
间冰期 $658.9 - 644.5 = \qty{14.4}{Ma}$。(b) $(58.0 + 9.1)/85 = 0.79$：约五分之四
的时长。(c) 其开始时间仅限定在 649.9 与 \qty{639.0}{Ma} 之间：按这些
界限，马里诺冰期的持续时间介于约 3 与 \qty{15}{Ma} 之间。
\end{solution}

\begin{solution}{exo:b2:precambrian-earth:8}
(a) $\tau = \frac{0.05}{0.95}\times\frac{44.01}{28.97}\times\frac{\num{5.1e18}}{\num{0.26e12}} =
0.05263\times1.519\times\num{1.962e7} = \qty{1.57}{Ma}$。(b) $0.05/\num{278e-6} = 180$ 倍。
(c) 时长加倍，为 \qty{3.14}{Ma}。
\end{solution}

\begin{solution}{exo:b2:precambrian-earth:9}
(a) 石块落入软泥，使其纹层向下弯曲：它们是落石，由
漂浮的冰释放。(b) 冰抵达距赤道 \ang{10} 的海面，是雪球态的
预期痕迹，若只有极地冻结则极难获得。(c) \omterm{def:b2:precambrian-earth:snowball}{盖帽碳酸盐岩}：
它记录了冰期的突然结束，伴随强烈风化，以及
融水泛滥于陆架之上。
\end{solution}

\begin{solution}{exo:b2:precambrian-earth:10}
(a) $0.94 > 0.932$，即暖态临界点：暖态平衡，$288\times0.94^{1/4} =
\qty{283.6}{K}$，略高于 \qty{283}{K}；且 $0.94 < 1.227$，故雪球态亦存在。
(b) $249.8\times0.94^{1/4} = \qty{246.0}{K}$。(c) $\varepsilon = 0.94\times1361\times0.4/(4\sigma
\times263^4) = 511.7/1085.1 = 0.472$。(d) $\sigma\times263^4 = \qty{271.3}{W.m^{-2}}$，且 $(0.616
- 0.472)\times271.3 = \qty{39}{W.m^{-2}}$.
\end{solution}

\begin{solution}{exo:b2:precambrian-earth:11}
(a) $S_{\text{m}}/S_{\text{f}} = (T_1/T_2)^4\,(1 - a_{\text{w}})/(1 - a_{\text{i}})$：
发射率相消。$(263/283)^4\times0.706/0.4 = 0.7459\times1.765 = 1.317$。(b) 冰与开阔水体之间反照率的
反差较小，例如含尘、较暗的冰，或更宽的
温度变化区间。(c) 雪球态反射的阳光如此之多，以至于在任何
阳光或温室效应下都保持寒冷（暖地球在此条件下会冻结）；要使其
融化，强迫量须超过冻结时的强迫量，超出幅度为整个回线的宽度。
\end{solution}

\begin{solution}{exo:b2:precambrian-earth:12}
在 $T^*$ 附近，$C\dot\delta = -(E' - R')\delta$，故 $\delta \propto e^{-t/\tau}$，其中 $\tau =
C/(E' - R')$。在暖态平衡处 $R' = 0$，而 $E' = 4\times0.616\times\num{5.670e-8}\times
288^3 = \qty{3.34}{W.m^{-2}.K^{-1}}$；$C = 1025\times3992\times50 = \qty{2.05e8}{J.m^{-2}.K^{-1}}$，
故 $\tau = \qty{6.1e7}{s}$，约 \qty{1.9}{\text{年}}。在本模型中，暖态平衡保持 $R' =
0$，直至达到 \qty{283}{K}，故 $\tau$ 在临界点之前约保持两年，随后
该态消失。若反照率平滑变化，则 $R'$ 随平衡态趋近
折叠点（$E' = R'$ 处）而增大：$\tau$ 无界增长，这种减慢预示着
临界点。
\end{solution}

\begin{solution}{pb:b2:precambrian-earth:1}
\textbf{1.} $\num{0.0203e-6}\times0.9929\times\num{95.13e-6} = \qty{1.92}{pW.kg^{-1}}$
(\ce{^{238}U}); $\num{0.0203e-6}\times0.00711\times\num{568.47e-6} = \qty{0.0820}{pW.kg^{-1}}$
(\ce{^{235}U}); $\num{0.0795e-6}\times\num{26.28e-6} = \qty{2.09}{pW.kg^{-1}}$ (\ce{^{232}Th});
$\num{240e-6}\times\num{1.20e-4}\times\num{28.47e-6} = \qty{0.820}{pW.kg^{-1}}$ (\ce{^{40}K}).
\textbf{2.} \qty{4.91}{pW.kg^{-1}}.
\textbf{3.} $\ln2/T$：0.155、0.985、0.0495 和 \qty{0.555}{Ga^{-1}}。
\textbf{4.} 因子为 1.52、14.3、1.14 和 4.48：2.91、1.17、2.39 和 \qty{3.67}{pW.kg^{-1}}。
\textbf{5.} \qty{10.1}{pW.kg^{-1}}，为今日的 2.07 倍。
\textbf{6.} 钾-40。其半衰期 \qty{1.25}{Ga} 较短：自那时起已过去两个多
半衰期，其时 \ce{^{40}K} 中有四分之三以上已衰变。
\textbf{7.} $1/(1 + 0.4\times2.7/4.567) = 0.809$; $0.809\times1361 = \qty{1101}{W.m^{-2}}$.
\textbf{8.} $288\times0.809^{1/4} = \qty{273}{K}$，低于 \qty{283}{K}：不是有效的暖态；
地球将会冻结。
\textbf{9.} $(1101\times0.4/(4\times0.616\,\sigma))^{1/4} = \qty{237}{K}$.
\textbf{10.} $0.616\times0.809 = 0.498$.
\textbf{11.} $0.004/\num{278e-6} = 14$ 至 $0.03/\num{278e-6} = 108$ 倍。
\textbf{12.} 硅酸盐风化恒温器：寒冷气候下风化减缓，使
火山排放的二氧化碳得以积累，而在温暖气候下则加快。
\textbf{13.} 为今日的 0.826，低于 0.932：不存在暖态；以今日的大气计，地球将成为
\omterm{def:b2:precambrian-earth:snowball}{雪球地球}。
\textbf{14.} $\frac{0.1}{0.9}\times1.519\times\num{1.962e7} = \qty{3.31}{Ma}$.
\textbf{15.} $\frac{0.2}{0.8}\times1.519\times\num{1.962e7} = \qty{7.45}{Ma}$.
\textbf{16.} \omterm{def:b2:precambrian-earth:snowball}{盖帽碳酸盐岩}，为一种层状白云岩，明显覆盖于冰川沉积之上，
它沉淀于由冰后炽热期强烈风化所补给的海水中，并
受融水泛滥。
\textbf{17.} 铁(II)仅在无氧水体中保持溶解：铁在
冰下缺氧的海洋中积累，在氧气到达处沉淀为铁(III)，例如
融水带来氧气之处。
\textbf{18.} $(1616 - 1450)/(32.2 - 23) = \qty{18.0}{\celsius}$ 每百分点。
\textbf{19.} $1616 + 18.0\times(30 - 32.2) = \qty{1576}{\celsius}$.
\textbf{20.} \qty{1616}{\celsius}，即实测点。
\textbf{21.} $1616 - 18.0\times20.2 = \qty{1252}{\celsius}$.
\textbf{22.} $1616 - 20\times20.2 = \qty{1212}{\celsius}$：低 \qty{40}{\celsius}，约为
熔岩之间差值的十分之一；结论并不依赖于这一选择。
\textbf{23.} \omterm{def:b2:precambrian-earth:greenstone}{科马提岩}来自地幔柱，其温度高于周围地幔，并源于深部的较大程度
熔融；熔岩的温度并不等于其源区的位温。整体而言，
地幔曾高出 200 至 \qty{300}{K}。
\textbf{24.} 最热的\omterm{def:b2:precambrian-earth:greenstone}{科马提岩}熔岩约高出 \textbf{\qty{364}{\celsius}}，比
今日洋脊的原生玄武岩。
\end{solution}
